\documentclass[planificacion.tex]{subfiles}
\begin{document}
\subsection*{\centering PROGRAMA \ano}

\begin{enumerate}
    \item PRINCIPIOS Y FENÓMENOS FLUIDODINÁMICOS 
    \begin{itemize}
        \item Reconocer e interpretar los
principios presentes en
fenómenos fluido-dinámicos.
        \begin{itemize}
            \item Descripción de los fenómenos hidrostáticos e
hidrodinámicos.
\item Análisis de los fluidos en reposo planteando la
relación entre fuerza y área como módulo introductorio
de Presión.
\item Experimentación e interpretación del Principio de
Pascal.
\item Experimentación e interpretación del Principio de
Arquímedes.
\item Aplicación del teorema fundamental de la Hidrostática
e identificación de los elementos de cálculo para
transmisiones hidráulicas.
\item Análisis de procesos hidrodinámicos mediante flujo de
los fluidos, caudal, ecuaciones de continuidad,
viscosidad, coeficientes, número específico.
\item Identificación de número de Reynolds utilizado para
identificar el tipo de flujo. 
Análisis de tipos de flujo en fluidos: régimen laminar y
turbulento.
\item Formulación del teorema de Bernoulli y sus
aplicaciones en pérdidas de carga y resistencia por el
frotamiento, fórmulas de pérdidas, dimensionamiento
de cañerías.
\item Comprobar a través de la experimentación el
fenómeno de cavitación, sus consecuencias y formas
de evitarlo.
\item Descripción de las ondas de choque generadas por un
objeto.
\item Observación de modelos y prototipos
hidrodinámicamente estables. 
        \end{itemize}
        \item Utilizar las leyes y sus fórmulas
en el cálculo de dispositivos
neumáticos y oleohidráulicos.
        \begin{itemize}
            \item  Interpretación y análisis de la simbología utilizada en cañerías, componentes e instalaciones relacionadas con fluidos.
    \item Explicitación de ecuaciones de hidrostática e hidrodinámica.
        \item Realización de cálculo de recipientes sometidos a
presión; ya sea en depósitos específicos, instalaciones
o estructuras depósitos.
      \end{itemize}
    
    \end{itemize}
    
\item INSTALACIONES Y EQUIPOS OLEOHIDRAULICOS E HIDRÁULICOS.
    \begin{itemize}
        \item Reconocer instalaciones y
equipos hidráulicos e
interpretar su funcionamiento.
    \begin{itemize}
         \item Identificación de equipos y componentes hidráulicos,
sus funciones principales y aplicaciones.
\item Caracterización de bombas hidráulicas: centrífugas y
bombas de desplazamiento positivo.
\item Identificación de los ensayos realizados a bombas
para su correcto funcionamiento y mejoras en su
rendimiento.
\item Identificación de circuitos de comando de
instalaciones, máquinas y componentes.
\item Reconocimiento de tipos de Turbinas: Pelton, Francis
Kaplan y tipo Bulbo
\item Descripción de principios y condiciones de
funcionamiento de equipos hidráulicos: bombas
radiales, bombas axiales, bombas mixtas, bombas de
vacío, impulsores axiales, radiales y mixtas.
\item Caracterización de las turbinas, su relación con el
rendimiento y las curvas de utilización.
\item Aplicación de los principios en frenos hidráulicos. 
    \end{itemize}
        \item Reconocer sistemas oleohidráulicos.
    \begin{itemize}
        \item Identificación de sistemas oleohidráulicos y circuitos
de presión: componentes básicos y lógica de
funcionamiento de los componentes.
\item Esquematización de los circuitos en sistemas oleohidráulicos y su aplicabilidad.
 Descripción de elementos y equipos oleohidráulicos:
prensas hidráulicas, sistemas de carga y transporte,
sistemas de elevación, sistemas para compactación.
\item Determinación de presiones admisibles en circuitos y
sistemas oleohidráulico e identificación de las medidas
de seguridad correspondientes.
    \end{itemize}
    \end{itemize}
    \item INSTALACIONES Y EQUIPOS NEUMÁTICOS. 
    \begin{itemize}
        \item Reconocer un sistema neumático.
        \begin{itemize}
            \item Identificación de los componentes de los sistemas neumáticos e interpretación de los principios teóricos necesarios.
            \item Discriminación de las variables intervinientes en el dimensionado de tuberías para el transporte de caudales.
            \item Reconocimiento de parámetros funcionales admisibles, presiones, velocidades, caudales.
            \item Interpretación y Descripción de las ecuaciones de fluidodinámica aplicables a los gases. 
        \end{itemize}
        \item Distinguir los sistemas de producción, de utilización y sus componentes.
        \begin{itemize}
            \item Clasificación en sistemas de producción y de utilización.
            \item Descripción de los componentes del sistema de producción: tipos de compresores, presostatos, cilindros, válvulas de seguridad, filtros, trampas de agua y lubricación.
            \item Cálculo de la capacidad en cilindros, presiones admisibles y métodos de protección ante la humedad; purgadores.
            \item Descripción de los componentes de un sistema de utilización neumático: tipos de válvulas, actuadores, elementos auxiliares, lógica de funcionamiento de los circuitos.
            \item Identificación y descripción de equipos y componentes electroneumáticos presentes en los subsistemas de producción y utilización que son posibles de automatizar.
            \item Identificación de ventiladores, soplantes y extractores y puesta a punto de dichos equipos. 
        \end{itemize}
    \item Automatizar sistemas neumáticos a través del uso de componentes. 
        \begin{itemize}
        \item Identificación de los componentes que pueden automatizarse.
        \item Análisis de las fuerzas necesarias para desplazar los accionamientos en los tiempos calculados.
        \item Aplicación de elementos en líneas de producción o sistemas en donde sea necesaria la automatización.
        \end{itemize}
    \end{itemize}
    \item ELEMENTOS DE MEDICIÓN Y ENSAYOS
    \begin{itemize}
        \item Reconocer los ensayos en instalaciones de fluidos. 
        \begin{itemize}
            \item Caracterización de tipos de ensayos referidos a fluidos.
            \item Realización de pruebas de estanqueidad y pruebas hidráulicas, de acuerdo a normas para la verificación y garantías en la calidad de la instalación. 
        \end{itemize}
        \item Realizar mediciones y reconocer instrumentos.
        \begin{itemize}
            \item  Determinación de mediciones de caudal y presión en cañerías y depósitos.
            \item Descripción de variables intervinientes en el cálculo y diseño de instalaciones hidráulicas.
            \item Identificación de los elementos de medición y ensayo: caudalímetro, manómetro relativo y absoluto, piezómetro. 
        \end{itemize}
    \end{itemize}
    
\end{enumerate}


\end{document}